%%%%%%%%%%%%%%%%%%%%%%%%%%%%%%%%%
%
%       EXERCÍCIO
%
%%%%%%%%%%%%%%%%%%%%%%%%%%%%%%%%%

\ifdefstring{\atividade}{prova}{%
    \renewcommand{\valorquestao}{\ValorQ\ pontos}
}%

\begin{exercicioBanco}[\valorquestao]
A tabela a seguir apresenta a distribuição de frequências dos salários (em salários mínimos) dos funcionários de uma empresa:

\begin{center}
\begin{tabular}{|c|c|}
\hline
Salário ($S_i$) & Frequência ($f_i$) \\ \hline
$2 \vdash 4$ & 4 \\ \hline
$4 \vdash 6$ & 10 \\ \hline
$6 \vdash 8$ & 4 \\ \hline
$8 \vdash 10$ & 2 \\ \hline
\end{tabular}
\end{center}

Com base nesses dados, determine:

\begin{enumerate}[(a)]
    \item A variância populacional ($\sigma^2$).
    \item O desvio padrão populacional ($\sigma$).
    \item O coeficiente de variação ($CV$).
\end{enumerate}
\end{exercicioBanco}
